本文围绕高阶 SDF 接触建模的适用域问题，提出并验证了一种\textbf{接触类型感知的高阶 SDF 接触框架}。与将单一接触代理统一施加于所有刚体接触场景的做法不同，本文将接触问题区分为光滑连续滑移接触与紧凑/奇异主导接触两类，并分别采用 \texttt{SlidingPatch} 与 \texttt{CompactDiscrete} 两种局部接触表示。二阶曲率项并不被视为一个对所有接触一概最优的统一算子，而是被放置在与接触类型相匹配的局部流形表示中发挥作用。

实验结果表明，这一建模思想在光滑连续接触中具有明确价值。对于工业凸轮、偏心圆盘-滚子解析基准以及 Onset Stress Benchmark，结合 persistent sliding manifold、局部几何拟合与面向 onset 的候选筛选的高阶 SDF 建模能够在保持显著效率优势的同时降低相对于一阶 SDF 的轨迹与速度误差，并在部分速度指标上逼近甚至略优于原生 mesh 接触。特别是在接触建立与持续滑移阶段，该框架能够抑制首次接触引发的局部峰值误差，并恢复后续的长期贴合行为。

与此同时，Twin Gear 基准也表明，高阶几何前馈并不能消除所有复杂接触中的误差来源。在带有 $C^0$ 奇异边界、最近点分支剧烈切换以及局部非光滑啮合的场景中，二阶 SDF 虽仍可取得相对于一阶 SDF 的净改善，但其误差下限仍然受到点接触微分展开对非光滑流形的先天敏感性所限制。因此，本文的核心结论并不是“统一二阶算法全面优于传统方法”，而是：\textbf{高阶 SDF 的有效性取决于接触类型及其局部流形表示，接触算法应按接触类型选择，而非由单一接触代理统治所有场景。}

未来的工作将沿两条路径继续推进。其一，是进一步完善面向连续接触的 persistent manifold 建模与局部机制消融，使接触建立、贴合保持与局部曲率修正之间的因果关系更加清晰。其二，是面向尖锐边缘及极小间隙主导的复杂装配体，推动接触建模由点导数范式（Point-based Differential Approximation）向体积积分 SDF（Volumetric Integral SDF）范式演进，以利用积分型几何测度作为天然低通滤波器，降低 $C^0$ 奇异点与分支切换带来的结构性高频误差。
